\section{Introduction}
\IEEEPARstart{N}{pm} dependency and maintainer relations can concentrate downstream reach in a small set of accounts. Earlier ecosystem studies measured this concentration and proposed maintainer vetting and metadata-based review signals~\cite{zimmermann2019,zahan2022}. A defensive program still needs to decide what a limited account-review budget should accomplish. Distinct repository coverage, exposure across account-origin scenarios, and selection of every observed source are different quantities.

Consider two accounts whose package portfolios reach many of the same repositories. Selecting the first may leave little additional coverage for the second. Nevertheless, an unselected second account remains a possible origin in the structural model. A repository can therefore be covered while retaining most of its observed sources. This is an objective distinction, not evidence that either account was compromised or that review would prevent compromise.

Our study extends established reach and greedy-coverage baselines~\cite{zimmermann2019,nemhauser1978} by measuring three objectives on the same observed incidence relation. We evaluate five deterministic selection rules, 200 random orders and six budgets. We then examine exact-version continuity, endpoint deduplication, observation exclusions, and hybrid and exclusion-oriented selections. The graph contains 25,228 identifiers and 12,217 usernames; workflow-positive package observations map to 7,171 repository keys. These are the units in a seeded expansion, rather than a representative sample of npm installations.

October 2026 reports illustrate boundaries that this graph cannot resolve. LiteLLM described master-key exposure, CI secret theft and misuse of developer tokens~\cite{litellm2026}, with published artifacts differing from a source tag~\cite{litellmincident2026}. The incident concerns PyPI/container releases, not an npm validation case. Pillar's OpenCode disclosure concerned a legacy npm dependency on Windows/WSL~\cite{opencode2026}. Railway described possible environment-variable exposure and a deployment pause~\cite{railway2026,railwaydisclosure2026}. SlugFest concerned publishing trust bindings~\cite{npmtrustedpub2026,npmexpiry2026}; the Node.js Brotli disclosure concerns a runtime outside the package graph~\cite{nodejsbrotli2026}. We retain these as contextual boundaries, without deriving intervention efficacy from them.

We ask five questions:
\begin{itemize}
\item \textbf{RQ1:} How do distinct coverage and additive scenario exposure differ at equal account budgets?
\item \textbf{RQ2:} How many covered repositories retain unselected sources, and what does an exclusion-oriented heuristic change?
\item \textbf{RQ3:} How do graph representation and selection/evaluation combinations affect the measured results?
\item \textbf{RQ4:} What tradeoffs arise under scalarized objectives and assumed heterogeneous weights?
\item \textbf{RQ5:} Which historical metadata can be reconstructed, and which evidence remains missing for authority-aware outcome evaluation?
\end{itemize}

The contribution is a measured comparison of established structural objectives, their representation and observation sensitivities, and the limits of the retained metadata. The scalarized objective has a standard submodular guarantee; our seven greedy settings do not certify a global Pareto frontier. The direct exclusion rule is a heuristic. The historical companion and frozen prospective acquisition pilot supply feasibility evidence, without a fitted forecasting model or deployed-control evaluation.
